%%%% necessary customization for the appendix and headers
\appendix 
\makeatletter
\addtocontents{toc}{\protect\renewcommand\protect\cftchappresnum{\@chapapp\ }}
\makeatother
\renewcommand{\thechapter}{\Alph{chapter}}

\chapter{Entropic Optimal Transport and the Sinkhorn Iteration} \label{chap:appendix-a}

%% Appendix A: the entropic-OT coupling + log-domain Sinkhorn used by StageBridge (Chapter 3, Eq. 3.9)
%% and as the endpoint/ladder distance. Matches crrt/ude/_ot.py: log-domain, float64, final rescaling.
%% VOICE: define every symbol; numbered equations; plain-word statement alongside the algebra.

StageBridge uses entropic optimal transport in two places: to build the source--target coupling
$\Pi^{(k)}$ that supervises training (Chapter~\ref{chap:methods}, Equation~\ref{eq:coupling}), and as the
Sinkhorn distance that measures endpoint accuracy in the falsification ladder
(Section~\ref{sec:ladder}). This appendix defines the entropic transport problem and the log-domain
Sinkhorn iteration used to solve it, matching the implementation.

\section{The entropic transport problem}

Given a source point set $\{z_i\}_{i=1}^{N}$ and a target point set $\{z_j\}_{j=1}^{M}$ in the
regulatory state space $\mathbb{R}^{d_z}$, define the cost matrix $C \in \mathbb{R}^{N \times M}$ by the
squared Euclidean distance,
\begin{equation}\label{eq:app-cost}
C_{ij} \;=\; \lVert z_i - z_j \rVert_2^2 ,
\end{equation}
where $C_{ij}$ is the cost of transporting a unit of mass from source $i$ to target $j$. The discrete
Kantorovich transport problem seeks a coupling $\Pi \in \mathbb{R}_{\geq 0}^{N \times M}$ that moves the
source distribution onto the target distribution at minimum total cost. With uniform marginals
$a = \tfrac{1}{N}\mathbf{1}_N$ and $b = \tfrac{1}{M}\mathbf{1}_M$, the admissible couplings are
\begin{equation}\label{eq:app-polytope}
\mathcal{U}(a, b) \;=\; \bigl\{\, \Pi \in \mathbb{R}_{\geq 0}^{N \times M} \;:\;
\Pi \mathbf{1}_M = a,\ \ \Pi^{\top} \mathbf{1}_N = b \,\bigr\},
\end{equation}
where $\mathbf{1}_N$ is the all-ones vector of length $N$, the row-sum constraint $\Pi \mathbf{1}_M = a$
fixes how much mass leaves each source point, and the column-sum constraint $\Pi^{\top}\mathbf{1}_N = b$
fixes how much arrives at each target point. The entropic problem adds a negative-entropy regularizer,
\begin{equation}\label{eq:app-entropic}
\Pi^{\star} \;=\; \argmin_{\Pi \in \mathcal{U}(a,b)}\ \sum_{i,j} \Pi_{ij}\, C_{ij}
\;-\; \epsilon\, H(\Pi),
\qquad
H(\Pi) \;=\; -\sum_{i,j} \Pi_{ij}\bigl(\log \Pi_{ij} - 1\bigr),
\end{equation}
where $H(\Pi)$ is the entropy of the coupling and $\epsilon > 0$ is the regularization weight
($\epsilon = 0.05$ throughout). The entropy term makes the objective strictly convex and the solution
unique and differentiable, which is what allows the coupling to be used inside a gradient-based training
loop; larger $\epsilon$ yields a smoother, more diffuse coupling, and $\epsilon \to 0$ recovers the
unregularized transport plan.

\section{The Sinkhorn fixed point}

The unique minimizer of Equation~\ref{eq:app-entropic} has the factored form
\begin{equation}\label{eq:app-factored}
\Pi^{\star}_{ij} \;=\; u_i\, K_{ij}\, v_j ,
\qquad
K_{ij} \;=\; \exp\!\left(-\,C_{ij} / \epsilon\right),
\end{equation}
where $K \in \mathbb{R}_{\geq 0}^{N \times M}$ is the Gibbs kernel of the cost and $u \in
\mathbb{R}_{\geq 0}^{N}$, $v \in \mathbb{R}_{\geq 0}^{M}$ are nonnegative scaling vectors. Imposing the
two marginal constraints of Equation~\ref{eq:app-polytope} on the factored form yields the Sinkhorn
fixed-point equations, whose alternating solution is the Sinkhorn iteration: repeatedly rescale $u$ so
the rows match $a$, then $v$ so the columns match $b$.

\section{Log-domain implementation}

Evaluating Equation~\ref{eq:app-factored} directly is numerically unstable because $K_{ij} =
\exp(-C_{ij}/\epsilon)$ underflows for the small $\epsilon$ used here. The implementation therefore
carries out the iteration entirely in the log domain, in double precision. Writing $\log K_{ij} =
-C_{ij}/\epsilon$ and log-scalings $\log u \in \mathbb{R}^{N}$, $\log v \in \mathbb{R}^{M}$, the updates
are
\begin{align}
\log u \;&\leftarrow\; \log a \;-\; \operatorname{logsumexp}_{j}\bigl(\log K_{ij} + \log v_j\bigr), \label{eq:app-logu}\\
\log v \;&\leftarrow\; \log b \;-\; \operatorname{logsumexp}_{i}\bigl(\log K_{ij} + \log u_i\bigr), \label{eq:app-logv}
\end{align}
where $\operatorname{logsumexp}(x) = \log \sum_k \exp(x_k)$ is evaluated by the standard max-shifted
form so no intermediate exponential overflows. Equations~\ref{eq:app-logu}--\ref{eq:app-logv} are
iterated for a fixed $80$ iterations; the coupling is then recovered as
\begin{equation}\label{eq:app-recover}
\log \Pi_{ij} \;=\; \log u_i \;+\; \log K_{ij} \;+\; \log v_j ,
\qquad
\Pi_{ij} \;=\; \exp\!\left(\log \Pi_{ij}\right),
\end{equation}
followed by a short sequence of alternating row/column rescalings of $\Pi$ that drives the marginals to
$a$ and $b$ to numerical tolerance. The entire computation is performed in double precision: single
precision was found to accumulate error across the thousands of in-loop gradient steps of training and
to drive the field parameters to non-finite values, whereas double precision is stable, and the
increased precision is worth the modest cost.

\section{Use as a distance}

The same routine supplies the endpoint distance of the falsification ladder. For a predicted endpoint
set $\{\hat z_i\}$ and a real target-stage set $\{z_j\}$, the Sinkhorn distance is the transport cost of
the optimal entropic coupling,
\begin{equation}\label{eq:app-distance}
d_{\mathrm{Sinkhorn}}\bigl(\{\hat z_i\}, \{z_j\}\bigr) \;=\; \sum_{i,j} \Pi^{\star}_{ij}\, \lVert \hat z_i - z_j \rVert_2^2 ,
\end{equation}
where $\Pi^{\star}$ is the coupling of Equation~\ref{eq:app-recover} computed between the two sets.
Equation~\ref{eq:app-distance} compares whole distributions of receivers rather than paired points,
which is the appropriate comparison when the prediction and the target are unpaired populations of
held-out cells.